%%% FIELD proof-main
Let \(\mathcal Q\) be the collection of quadruples
\[
 Q=(k,R_1,R_2,W),
 \qquad k\in[K],
 \qquad R_1,R_2,W\subseteq U_k,
\]
for which \(R_1,R_2,W\) are pairwise disjoint and \(R_1\) and \(R_2\) are \(m\)-connected given \(W\) in the replacement MAG on \(U_k\).  This collection is finite: \([K]\) and every \(U_k\) are finite.

Fix \(Q=(k,R_1,R_2,W)\in\mathcal Q\) and a joint cell assignment \(a=(r_1,r_2,w)\).  Define
\[
 \Delta_{Q,a}(\theta)
 :=p_{\theta,k}(r_1,r_2,w)p_{\theta,k}(w)
   -p_{\theta,k}(r_1,w)p_{\theta,k}(r_2,w).
 \tag{1}
\]
Here each displayed probability is the appropriate marginal of the pooled mass function from the same coupled-response model \(\eta_\theta\).  Strict positivity of every full cell implies \(p_{\theta,k}(w)>0\) for every \(\theta\in\Theta\) and every \(w\); if \(W=\varnothing\), this marginal equals \(1\).  Consequently, the finite-table definition of conditional independence gives
\[
 X_{R_1}\mathbin{\perp\!\!\!\perp}X_{R_2}\mid X_W
 \quad\Longleftrightarrow\quad
 p_{\theta,k}(r_1,r_2,w)p_{\theta,k}(w)
 =p_{\theta,k}(r_1,w)p_{\theta,k}(r_2,w)
 \quad\text{for every }a.
 \tag{2}
\]
Indeed, division by \(p_{\theta,k}(w)^2\) turns the equality on the right into the usual factorization of the conditional joint mass, and positivity justifies that division.  Thus (1)--(2) imply
\[
 X_{R_1}\mathbin{\perp\!\!\!\perp}X_{R_2}\mid X_W
 \quad\Longleftrightarrow\quad
 \Delta_{Q,a}(\theta)=0\quad\text{for every }a.
 \tag{3}
\]

All full-cell masses are rational in \(\theta\), so their finite sums, and hence all marginals in (1), are rational in \(\theta\).  By the stipulated positivity of all denominators on \(\Theta\), clearing denominators in (1) yields the polynomial named in the statement and a function \(d_{Q,a}:\Theta\to(0,\infty)\) such that
\[
 q_{Q,a}(\theta)=d_{Q,a}(\theta)\Delta_{Q,a}(\theta).
 \tag{4}
\]
It follows that, on \(\Theta\),
\[
 q_{Q,a}(\theta)=0
 \quad\Longleftrightarrow\quad
 \Delta_{Q,a}(\theta)=0.
 \tag{5}
\]

Assume first that for every \(Q\in\mathcal Q\) there is at least one assignment \(a(Q)\) for which \(q_{Q,a(Q)}\) is not the zero polynomial.  Set
\[
 q_*(\theta):=\prod_{Q\in\mathcal Q}q_{Q,a(Q)}(\theta),
 \tag{6}
\]
where the empty product is \(1\).  The ring \(\mathbb R[\theta_1,\ldots,\theta_r]\) is an integral domain, so finiteness of \(\mathcal Q\) makes \(q_*\) a nonzero polynomial.

We prove directly the elementary zero-set fact needed here.  A nonzero real polynomial in \(r\) variables cannot vanish on a nonempty open subset of \(\mathbb R^r\).  For \(r=0\) this says that a nonzero constant does not vanish.  For \(r=1\), a nonzero polynomial has at most its degree many distinct roots, by repeated division by \(t-t_0\).  Inductively, write
\[
 h(t_1,\ldots,t_r)=\sum_{j=0}^{d}h_j(t_1,\ldots,t_{r-1})t_r^j.
 \tag{7}
\]
If \(h\) vanished on an open box, then at every point in the box's projection the resulting univariate polynomial in \(t_r\) would vanish on an interval.  Every coefficient \(h_j\) would therefore vanish throughout the projected open box, and the induction hypothesis would make every \(h_j\) the zero polynomial, contradicting \(h\ne0\).  Since every nonempty open subset of Euclidean space contains an open box, the fact follows.

Apply this fact to the nonzero polynomial \(q_*\) and the nonempty open set \(\Theta\).  There is \(\theta^\star\in\Theta\) such that
\[
 q_*(\theta^\star)\ne0.
 \tag{8}
\]
For every \(Q\in\mathcal Q\), equations (6) and (8) give \(q_{Q,a(Q)}(\theta^\star)\ne0\).  Equations (3)--(5) then yield
\[
 X_{R_1}\not\!\mathbin{\perp\!\!\!\perp}X_{R_2}\mid X_W
 \quad\text{under }P_{\eta_{\theta^\star},k}^{\mathrm{pool}}
 \qquad\text{for every }Q\in\mathcal Q.
 \tag{9}
\]
For every query that is \(m\)-separated in its replacement MAG, the Markov-direction premise gives the corresponding conditional independence.  A query not \(m\)-separated is, by definition, one of the \(m\)-connected queries covered by (9).  Hence, for every \(k\in[K]\) and every pairwise disjoint \(R_1,R_2,W\subseteq U_k\),
\[
 X_{R_1}\mathbin{\perp\!\!\!\perp}X_{R_2}\mid X_W
 \quad\Longleftrightarrow\quad
 R_1\text{ and }R_2\text{ are }m\text{-separated by }W
 \text{ in the replacement MAG}.
 \tag{10}
\]
Thus all pooled selected laws are simultaneously faithful to their replacement MAGs.

By hypothesis, the directly parameterized tuple \(\eta_{\theta^\star}\) already satisfies every other member condition of \(\mathrm{def:model\mbox{-}class}\): in particular it has one product exogenous law with the same realized exogenous vector in every world, shared response functions at non-target vertices, one pre-intervention basal selection response, and positive basal selection probability.  Equation (10) supplies the only condition that was allowed to be absent, namely pooled twin faithfulness.  Therefore
\[
 \eta_{\theta^\star}\in\mathfrak M_{D,K}^{\mathrm{Dai}}.
 \tag{11}
\]
The family premise also says that the selection-set, basal-graph, and target components of \(\eta_{\theta^\star}\) jointly equal \(\operatorname{RepGraph}(\eta,B)\).  Applying \(\mathrm{def:replacement\mbox{-}realization\mbox{-}fiber}\) to (11) consequently gives
\[
 \eta_{\theta^\star}\in\operatorname{Real}(\eta,B).
 \tag{12}
\]
The cited frozen result \(\mathrm{lem:graphical\mbox{-}replacement\mbox{-}preservation}\) guarantees that this replacement specification has the asserted selected-DAG and regime-signature properties.  It is not used to construct or recouple the shared-response law; that law is supplied by the parameter-family premise.

Conversely, suppose that some \(\theta^\star\in\Theta\) makes every pooled selected law faithful to its replacement MAG.  Fix \(Q\in\mathcal Q\).  Because this query is \(m\)-connected, faithfulness implies
\[
 X_{R_1}\not\!\mathbin{\perp\!\!\!\perp}X_{R_2}\mid X_W.
\]
By (3), there is a cell assignment \(a\) for which \(\Delta_{Q,a}(\theta^\star)\ne0\).  Equations (4)--(5) imply
\[
 q_{Q,a}(\theta^\star)\ne0,
 \tag{13}
\]
and therefore \(q_{Q,a}\) is not the zero polynomial.  Since \(Q\) was arbitrary, the required atomic nontriviality holds for every \(m\)-connected query.

Finally, if \(K=0\), there are no pooled laws and \(\mathcal Q=\varnothing\).  The atomic condition and simultaneous-faithfulness condition are both vacuous, (6) gives \(q_*=1\), and any point of the nonempty set \(\Theta\) satisfies all member conditions because pooled twin faithfulness has no instances.  The same component-equality argument then gives (12).  This proves both directions, including the boundary case.
